\scanpage{74}
seconde, et je dis que ce foncteur est une équivalence de catégories.

Le cas qui m'intéresse est celui où $\underline G=\mu_2\unsure{\simeq|\xrightarrow{\sim}|\underset{\mathrm{can}}{\simeq}}\{\pm1\}_S$, et où $n=2$ -- donc où on s'intéresse aux rev. quadratiques $X'$ de $X/S$, étales au dessus de $X$ et ``totalement ramifiés le long de la section $t$'', avec $S$ à car. rés $\ne 2$. \deleted{\unsure{Si|Soit|Ici}} Pour une section $\xi\in H^0(S,R^1f_\circ^\ast(\mu_{2\,X^\ast}))$ fixée, satisfaisant à la condition de ramification totale, la catégorie des \deleted{\unsure{ce|ces|cou}}revêtements \inserted{correspondant à cette $\xi$} est donc équivalente à celle des \deleted{\unsure{rev.|rev}} couples $(T,k)$ d'un Module inv. $T$ sur $S$, et d'un isom
\[
T^{\otimes2}\simeq T_{X,t} .
\]

Le cas qui nous intéresse pour les courbes elliptiques est celui où $X$ est de la forme $\Sigma_{J}^\ast$ (fibré en droites projectives associé à un rev. étale cubique $J$ de $S$), \inserted{déduit d'une courbe elliptique $X'$ comme $X'/\{\pm1\}$,} et où $t$ est la section de $\Sigma_J^\ast/S$ \deleted{\unsure{quotient|qui}} image de la section nulle de $X'=E$ \unsure{en|un}
% archive sheet number: 73
